\documentclass[a4paper,10pt]{article}
\usepackage[utf8]{inputenc}
\usepackage[T1]{fontenc}
\usepackage{parskip}
\usepackage[sfdefault]{FiraSans}
\usepackage[table]{xcolor}
\usepackage{colortbl}
\usepackage{geometry}
\geometry{margin=2cm}
\usepackage{array}
\usepackage{enumitem}
\usepackage{titlesec}
\usepackage{hyperref}
\hypersetup{colorlinks=false}

\definecolor{ci-dark}{HTML}{2D3748}
\definecolor{ci-accent}{HTML}{2B6CB0}
\definecolor{ci-light}{HTML}{718096}
\definecolor{ci-rule}{HTML}{CBD5E0}
\definecolor{ci-bg}{HTML}{EDF2F7}
\definecolor{ci-total}{HTML}{1A202C}

% Fix titlesec: replace \\ before-separator with empty to prevent "0em[]" artifact
\titleformat{\section}{\large\bfseries\color{ci-dark}}{}{0em}{\vspace{0.1em}}
\titleformat{\subsection}{\normalsize\bfseries\color{ci-accent}}{}{0em}{}

\pagestyle{empty}
\setlength{\parindent}{0pt}
\setlength{\parskip}{2pt plus 0.5pt minus 0pt}
\setlength{\arrayrulewidth}{0.4pt}
\setlist{nosep,leftmargin=1.5em}

\begin{document}

% === HEADER ===
\begin{center}
\vspace{0.2cm}
{\LARGE\bfseries\color{ci-dark} Geschäftsmodellskizze}\par
\vspace{0.2cm}
{\large\color{ci-accent} StartMiUp Businessplan-Wettbewerb 2026}\par
\vspace{0.5cm}

\begin{tabular}{@{}l @{\hspace{1em}}l}
\textcolor{ci-light}{\textbf{Projekt:}}       & \textbf{Contextual Intelligence (CI)} \\
                                      & \textcolor{ci-light}{KI-gestützte Lead- \& Grant-Intelligence-Plattform} \\[0.35em]
\textcolor{ci-light}{\textbf{Antragsteller:}} & Tobias \\
                                      & tobias@tobias-weiss.org \textcolor{ci-rule}{\textbullet}~ 0177-9132333 \\[0.35em]
\textcolor{ci-light}{\textbf{Hochschule:}}   & Philipps-Universität Marburg, HRZ \\[0.35em]
\textcolor{ci-light}{\textbf{Datum:}}         & 26.06.2026 \\
\end{tabular}

\vspace{0.3cm}
\textcolor{ci-accent}{\rule{\textwidth}{0.6pt}}
\end{center}

\vspace{0.15cm}

% === SECTION 1 ===
\section{Einleitung}

Contextual Intelligence (CI) ist eine KI-SaaS-Plattform für B2B-Lead-Intelligence im Deep-Tech-Bereich. Ein semantischer Wissensgraph auf selbst-gehosteten Foundation-Modellen identifiziert Technologie-Partner automatisch — statt starrer Keywords. Das MVP (Lead Intelligence, Siemens-Hackathon) ist live.

\textbf{Problem:} Deep-Tech-Vertrieb ist intransparent; bestehende Tools (Apollo, ZoomInfo, LinkedIn Sales Navigator) haben null Technologie-Kontext und sind keyword-basiert.

\textbf{Relevanz:} Bessere Leads = mehr Umsatz bei weniger Aufwand.

% === SECTION 2 ===
\section{Wertangebot}

\textbf{Zielgruppe:} Deep-Tech-B2B-Vertriebsteams (Hardware, Robotics, Bio, AI), Forschungsorganisationen, F\&E-KMU.

\textbf{Problem:} Bestehende B2B-Datenbanken (Apollo, ZoomInfo, LinkedIn Sales Navigator) sind generisch und keyword-basiert — sie verstehen keine spezialisierten Technologien und deren Überlappungen. Cold Outreach ohne Signal-Bezug hat niedrige Conversion. Zusätzlich: 1.500+ Förderprogramme in DE/EU werden manuell recherchiert.

\textbf{USP:} CI kombiniert:
\begin{enumerate}[nosep,leftmargin=1.8em]
  \item \textbf{Lead Intelligence} (live) + \textbf{Grant Intelligence} (geplant)
  \item \textbf{Semantischer Wissensgraph} (Dgraph) — kontextuelles Verständnis statt Keywords
  \item \textbf{Sovereign-AI-Infrastruktur} — 2$\times$ NVIDIA DGX Spark, Open-Weight-Foundation-Modelle, vollständig DSGVO-/EU-AI-Act-konform, kein Datenabfluss in Drittstaaten
\end{enumerate}
Vorteile: Graph-basiertes Scoring (HOT/WARM/COLD), Fixkostendominiert (Cost-per-Scan <0,001 Euro) statt pro-Token-API-Kosten.

% === SECTION 3 ===
\section{Geschäftsmodell-Logik (Business Model Canvas)}

\begin{description}[nosep,leftmargin=3em,style=nextline,itemsep=1.5pt]
  \item[Kanäle] Product-led Growth — Self-Service (Website, Stripe, API-Key), Content-Marketing, Open-Source (ORSD als Lead Magnet), Direktkontakt nur bei Enterprise.
  \item[Kundenbeziehungen] Self-Service First (Docs, FAQ, Community); persönlicher Support nur ab €1.000/Monat. Vision: weitgehend autonome Business Unit — das Produkt verkauft und onboardet sich selbst.
  \item[Einnahmequellen] Lead Intelligence €299/Monat (B2B-Vertriebsteams); Lead Enterprise ab €1.000/Monat; Grant Intelligence €49–€149/Monat (geplant Q2 2027); ORSD Premium-API €99–€499/Monat. Zusätzlich: Forschungszulage (BZSt, 25\,\% F\&E-Personal), Förderprojekte (LOEWE, BMBF).
  \item[Schlüsselressourcen] 2$\times$ NVIDIA DGX Spark (128\,GB, eingelegt), Dgraph-Wissensgraph, Postgres, React 19 + Express + Nuxt 3 (Open Source), ORSD-Dataset (6 Tabellen, 22 Crawler), z.ai-API als Fallback.
  \item[Schlüsselaktivitäten] Foundation-Model-Betrieb, Crawler-Wartung, Scoring-Tuning, Content-Marketing, ORSD-Community-Aufbau.
  \item[Schlüsselpartner] Hetzner (Hosting), z.ai (API), Uni Gießen (LOEWE-Verbund), Datenquellen (FDA, EUDAMED, Förderdatenbanken), StartMiUp/Hochschulnetzwerk (Pilotkunden).
\end{description}

\vspace{0.3em}
\nopagebreak
\textbf{Kostenstruktur (Y1, monatlich):}

\vspace{0.25em}

\begin{center}
\renewcommand{\arraystretch}{1.25}
\setlength{\tabcolsep}{8pt}
\begin{tabular}{>{\raggedright\arraybackslash}p{6cm} >{\raggedleft\arraybackslash}p{3cm}}
\rowcolor{ci-bg}
\textbf{\color{ci-dark}Posten} & \textbf{\color{ci-dark}Kosten} \\[2pt]
GPU-Strom + AfA            & 220 Euro \\[1pt]
z.ai API                    & 50--100 Euro \\[1pt]
Hetzner VPS + Docker         & 150--300 Euro \\[1pt]
Domain/SSL/Tools             & 20--50 Euro \\[1pt]
Freelancer                  & 300--600 Euro \\[1pt]
Marketing/Vertrieb            & 200--400 Euro \\[1pt]
Recht/Steuer/Admin           & 100--200 Euro \\[2pt]
\rowcolor{ci-bg}
\textbf{\color{ci-total}Gesamt-Burn} & \textbf{\color{ci-total}1.000--2.000 Euro/Monat} \\
\end{tabular}
\end{center}

\vspace{0.3em}
Vorteil Nebenerwerb: Kein Gründungsgehalt; Break-Even bei <2.000 Euro MRR (Monat 12--18 erreichbar).

\nopagebreak
% === SECTION 4 ===
\section{Kreativität \& Umsetzbarkeit}

\textbf{Innovation — drei USPs im DACH-Markt:}
\begin{enumerate}[nosep,leftmargin=1.8em]
  \item \textbf{Sovereign AI} — Open-Weight-Modelle auf DGX (kein US-Hyperscaler)
  \item \textbf{Graph-basiertes Lead-Scoring} über Technologie-Grenzen
  \item \textbf{Product-led Growth} ohne Sales-Team — das Produkt verkauft sich selbst
\end{enumerate}

\textbf{Umsetzbarkeit:} Lead-Intelligence-MVP live (Siemens-Hackathon). Website produktiv, 2$\times$ DGX Spark vorhanden. Y1: 10--20 Pilotkunden (€9k--€18k ARR). Funding: ERP-Gründerkredit €125k + LOEWE €100k--500k + BMBF.

\end{document}
